\documentclass{article}

% Load your code directly from the working directory:
\usepackage{pxPRP}

\begin{document}

\section*{Proves bàsiques de pxPRP}


% --- Prova 1: smoke test, creació i accés bàsic --------------------------------

\subsection*{Prova 1: Creació i accés bàsic}
\verb*|\pxNewObject{PATATA}|\\
\pxNewObject{PATATA}

\verb*|\pxSet{PATATA}{k1}{v1,v2}|\\
\pxSet{PATATA}{k1}{v1,v2}

\verb*|\pxSet{PATATA}{k2}{vA,vB,vC}|\\
\pxSet{PATATA}{k2}{vA,vB,vC}

\verb*|\pxGet PATATA.k1[2]| $\to$ \pxGet PATATA.k1[2]\\

\verb*|\pxPrint PATATA.k2[2]|$\to$ \pxPrint PATATA.k2[2]

\bigskip

\subsection*{Prova 2: Múltiples objectes i claus aïllades}

\pxNewObject{OBJAA}
\pxNewObject{OBJBB}

\pxSet{OBJAA}{key}{AA1,AA2}
\pxSet{OBJBB}{key}{BB1,BB2}

En pantalla: OBJAA.key[1] = \pxGet OBJAA.key[1],\quad
OBJBB.key[1] = \pxGet OBJBB.key[1].

\bigskip

\subsection*{Prova 3: Sobreescriptura amb \texttt{\textbackslash pxSet}}

\pxNewObject{OVER}
\pxSet{OVER}{state}{old}

Abans de sobreescriure: OVER.state[1] = \pxGet OVER.state[1].

\pxSet{OVER}{state}{newA,newB}

Després de sobreescriure: OVER.state[1] = \pxGet OVER.state[1],\quad
OVER.state[2] = \pxGet OVER.state[2].

\bigskip

\subsection*{Prova 4: Acumulació de llistes amb \texttt{\textbackslash pxPutValue}}

\pxNewObject{LST}

\pxPutValue{LST}{items}{a}
\pxPutValue{LST}{items}{b}
\pxPutValue{LST}{items}{c}

Resultat esperat: LST.items = (\pxGet LST.items[1], \pxGet LST.items[2], \pxGet LST.items[3]).

\bigskip
%
\subsection*{Prova 5: Esborrar claus i objectes}

\pxNewObject{TEMP}
\pxSet{TEMP}{k}{X,Y}

Abans d'esborrar: TEMP.k[1] = \pxGet TEMP.k[1].

\pxClearKey{TEMP}{k}

Després d'esborrar la clau, podem tornar a definir-la:

\pxSet{TEMP}{k}{Z}

Ara: TEMP.k[1] = \pxGet TEMP.k[1].

\pxClearObject{TEMP}
%
%Després d'esborrar l'objecte, podem recrear-lo sense errors:
%
%\pxNewObject{TEMP}
%\pxSet{TEMP}{k}{NOU}
%%
%TEMP.k[1] = \pxGet TEMP.k[1].
%%
\bigskip

\subsection*{Prova 6: Flags booleans}

\pxNewObject{FLAGS}

\pxSet{FLAGS}{featureA}{true}
\pxSet{FLAGS}{featureB}{{false}}

\pxIfTrueTF{FLAGS}{featureA}{ACTIVA}{INACTIVA}

\par

\pxIfTrueTF{FLAGS}{featureB}
  {featureB és ACTIVA.}
  {featureB és INACTIVA.}

\bigskip
%
%\bigskip
%\subsection*{Prova 8: Llistes amb \texttt{\textbackslash pxListPutLeft}}
%
%\pxNewObject{OBJL}
%
%\pxListPutRight{OBJL}{nums}{2,3}
%\pxListPutLeft{OBJL}{nums}{1}
%
%Esperat: OBJL.nums = (\pxGet OBJL.nums[1], \pxGet OBJL.nums[2], \pxGet OBJL.nums[3]).
%
%\bigskip
%
%\subsection*{Prova 9: Substitució completa amb \texttt{\textbackslash pxListSet}}
%
%\pxNewObject{OBJS}
%
%\pxListPutRight{OBJS}{vals}{a,b,c}
%
%Abans de \pxListSet: OBJS.vals[1] = \pxGet OBJS.vals[1].
%
%\pxListSet{OBJS}{vals}{x,y}
%
%Després de \pxListSet: OBJS.vals[1] = \pxGet OBJS.vals[1],\quad
%OBJS.vals[2] = \pxGet OBJS.vals[2].

\end{document}